\chapter{Profiling, performance and hardware backends}

Performance engineering begins with a refusal: do not optimize what you have not measured, and do not promote a microbenchmark into a system claim.

For Evolutor the measurement ladder is:
\[
\text{SLO}\rightarrow\text{request}\rightarrow\text{phase}\rightarrow
\text{operator}\rightarrow\text{kernel}\rightarrow\text{hardware counter}.
\]
\Plate{01-ladder}{Profile top-down from user-visible behavior.}{evod47:ladder}

\section{Correctness before speed}
Every optimization is paired with a correctness oracle. Compare logits/state/output distributions as appropriate before comparing time. A 2x speedup that changes semantics outside the declared numerical tolerance is a different implementation, not a valid optimization.

\section{Workload contract}
Record prompt/input length distribution, output length, concurrency/arrival process, model/artifact, precision/quantization, sampling, cache state, scheduler, batch limits and memory budget.

Performance without workload definition is not reproducible.

\section{Latency decomposition}
For one request:
\[
L_{e2e}=L_q+L_{prefill}+L_{transfer}+L_{decode}+L_{tool/memory}+L_{serialization}.
\]
Terms depend on architecture. DOGMA may have native state-ingestion/step phases rather than Transformer prefill/KV semantics.

\section{TTFT and ITL}
Time to first token:
\[
TTFT=t_{first}-t_{arrival}.
\]
Inter-token latency:
\[
ITL_i=t_i-t_{i-1}.
\]
\Plate{02-metrics}{TTFT and ITL expose different serving phases.}{evod47:metrics}

For Transformer systems, TTFT is strongly influenced by queue/prefill and ITL by decode scheduling. Report distributions, not a single average.

\section{Throughput and goodput}
Raw throughput may be output tokens/s or requests/s. Goodput counts only work satisfying the declared SLO:
\[
G=\frac{\#\{r:\ r\text{ successful and SLO-compliant}\}}{\Delta t}.
\]
A system can increase raw throughput while making most interactive requests unusable.

\section{Percentiles}
For a benchmark population report at least p50 and a tail percentile such as p95/p99. State the percentile estimator for small samples. Tail latency often reflects queueing, synchronization and interference rather than one slow kernel.

\section{Warmup and steady state}
Separate cold artifact loading, compilation/graph capture, cache warming and allocator initialization from steady-state execution unless cold start is itself the metric.

Use explicit warmup iterations and record them.

\section{Repetitions and uncertainty}
Repeat experiments across runs/seeds/workload traces. Report sample count and dispersion/confidence intervals where appropriate. Do not report more precision than the measurement supports.

\section{Arithmetic intensity}
For an operator with $F$ floating-point operations and $B$ bytes transferred from the limiting memory level:
\[
I=\frac{F}{B}\quad[\mathrm{FLOP/byte}].
\]
\Plate{03-roofline}{Arithmetic intensity connects algorithm work to hardware ceilings.}{evod47:roofline}

The relevant byte count depends on cache reuse and memory hierarchy; naive tensor-size sums may misestimate actual traffic.

\section{Roofline}
Given peak compute $P_{peak}$ and bandwidth $BW$:
\[
P\le \min(P_{peak}, BW\cdot I).
\]
The ridge point is
\[
I^*=\frac{P_{peak}}{BW}.
\]
If $I<I^*$, the roofline region is memory-bandwidth limited; otherwise the compute ceiling dominates. This is an upper-bound model, not achieved performance.

\section{Attention example}
Decode attention often reads a growing KV history for relatively little new-token work, making memory traffic central. FlashAttention-style tiled algorithms change IO behavior by increasing on-chip reuse; Chapter 38 derived online softmax correctness. Chapter 47 asks whether the intended IO reduction occurs on the target hardware.

\section{DOGMA example}
A recurrent/native DOGMA step may touch a bounded state representation each token. Its arithmetic intensity depends on actual state width, gates, memory banks, precision and kernel fusion. Do not infer performance from ``linear recurrence'' alone.

\section{System timeline}
\Plate{04-timeline}{Timeline profiling exposes gaps between useful kernels.}{evod47:timeline}
Measure CPU scheduling, kernel launches, synchronization, host/device copies, collective communication, allocator stalls and idle gaps.

NVIDIA recommends Nsight Systems for system-wide CPU/GPU tracing and Nsight Compute for detailed CUDA kernel analysis. citeturn0search1turn0search0

\section{Focused profiling}
Profile representative critical regions rather than drowning in initialization. Nsight Systems supports focused regions/NVTX annotations; this lets Evolutor label prefill, decode, DOGMA step, retrieval, transfer and tool phases on one timeline. citeturn0search6

\section{Operator profiling}
PyTorch Profiler can identify expensive model operators, inspect input shapes/stacks, device kernel activity and execution traces. This is useful for training and inference prototypes before descending into backend-specific kernels. citeturn0search5

\section{Kernel profiling}
Once a kernel is proven to lie on the critical path, inspect achieved occupancy, memory throughput, instruction mix, cache behavior, stalls and launch configuration with the backend's kernel profiler.

Counter interpretation is architecture-specific; collecting many counters can perturb execution.

\section{CPU backend}
CPU profiling emphasizes vectorization, cache locality, NUMA placement, thread pools, memory bandwidth and branch/dispatch overhead. Small-batch/edge workloads can be dominated by launch/framework overhead rather than nominal FLOPs.

Record ISA, core count, affinity and thread settings.

\section{Metal backend}
On Apple hardware, record exact SoC, macOS, Metal/runtime/framework version, unified-memory pressure and precision. Unified memory changes transfer mechanics but does not eliminate bandwidth/capacity limits.

Use platform-native GPU/system profiling for real claims; do not translate CUDA counters mechanically into Metal terminology.

\section{CUDA backend}
CUDA offers mature system and kernel profiling through Nsight Systems/Compute. Record GPU model, driver/toolkit, clocks/power mode, MIG/MPS state, interconnect and library/kernel versions.

Peak Tensor Core specifications are not achieved application throughput.

\section{Vulkan backend}
Vulkan compute portability can broaden hardware coverage. Performance depends on driver, shader compiler, subgroup features, memory layout and synchronization. Report exact device/driver and use platform/vendor profiling tools where available.

A portable API does not imply portable performance.

\section{Backend comparison}
\Plate{05-backends}{Compare backends only under semantic and workload parity.}{evod47:backends}
For CPU/Metal/CUDA/Vulkan compare:
\begin{enumerate}
\item identical model/state semantics and numerical tolerance;
\item identical prompt/output workload;
\item comparable precision/quantization;
\item warmup and steady-state protocol;
\item latency distributions and throughput/goodput;
\item peak memory and energy/power when measured;
\item backend-specific profile explaining the bottleneck.
\end{enumerate}

\section{Batching}
Batching can raise device utilization while increasing queue delay or ITL. Plot performance against concurrency/batch occupancy instead of reporting one ``optimal'' batch in isolation.

Continuous batching should be evaluated under an arrival process, not only a preconstructed full batch.

\section{Synchronization}
A synchronization inserted for timing can itself change execution. GPU timing requires device-aware events or profiler timelines. Wall-clock timing must synchronize at the measurement boundary intentionally and consistently.

\section{Distributed performance}
For Chapter 46 placement:
\[
L_{migration}=L_{queue}+L_{serialize}+L_{network}+L_{deserialize}+L_{commit}.
\]
Measure each term. Prefill/decode disaggregation can reduce phase interference, as DistServe reports, but introduces transfer/topology costs and must be evaluated under TTFT/TPOT constraints. citeturn0academia36

\section{Speedup arithmetic}
Speedup is
\[
S=\frac{T_{baseline}}{T_{new}}.
\]
Reducing 100 ms to 80 ms is $1.25\times$ speedup, not ``20x''. Percent latency reduction is a different statistic:
\[
\frac{100-80}{100}=20\%.
\]

\section{Benchmark record}
\Plate{06-record}{A performance claim includes enough context to reproduce it.}{evod47:record}
Record commit/artifact hashes, hardware, firmware/driver, OS, compiler/runtime, backend libraries, power/clocks, precision, scheduler, workload generator, warmup, repetitions, raw samples and analysis script.

\section{Executable analysis oracle}
\Listing{reference}
The reference computes percentile latency, TTFT/ITL, token throughput, SLO-qualified goodput, arithmetic intensity, roofline bounds and speedup from synthetic inputs.

Its numbers are deliberately synthetic and must never appear as measured hardware results.

\section{Optimization loop}
Use:
\[
measure\rightarrow locate\ bottleneck\rightarrow hypothesize\rightarrow
change\ one\ mechanism\rightarrow correctness\ test\rightarrow remeasure.
\]
Keep failed optimizations in the experiment ledger; they prevent repeated folklore.

\section{Twelve worked exercises}
\begin{enumerate}
\item Fastest kernel means fastest service? \textbf{No; it must be on the critical path and system effects matter.}
\item TTFT versus ITL? \textbf{First-token latency versus spacing between generated tokens.}
\item Throughput versus goodput? \textbf{Goodput counts only SLO-compliant successful work.}
\item $I=2$ FLOP/B, $BW=2$ TB/s: bandwidth roof? \textbf{4 TFLOP/s.}
\item Roofline bound equals achieved throughput? \textbf{No.}
\item Why warmup? \textbf{Separate initialization/compilation/cache effects from steady state.}
\item Why profile system before kernel? \textbf{The bottleneck may be queueing, copies, launch or synchronization.}
\item CUDA peak spec enough? \textbf{No.}
\item Unified memory means unlimited/free memory traffic? \textbf{No.}
\item 20 ms to 10 ms speedup? \textbf{2x.}
\item Can synthetic oracle results be quoted as GPU benchmarks? \textbf{No.}
\item Release rubric? \textbf{Correctness parity; workload contract; percentiles; TTFT/ITL/goodput; warmup/repeats; top-down profiles; roofline with valid bytes; backend versions; raw evidence; reproducible analysis.}
\end{enumerate}

\section{Bridge}
The industrial runtime foundation is now complete enough to return to applications. Chapter 48 applies the model/runtime stack to genomic sequence tasks while preserving the book's central discipline: biological sequence is data with domain-specific structure, not proof that a computational mechanism is biologically faithful.
